\section{Introduction}
\label{sec:introduction}

Memory-augmented Large Language Models (LLMs) are increasingly deployed in long-running dialogue settings, but most current systems operate under what is essentially a FIFO or accumulation-only regime: content is added to a vector store and retrieved by semantic similarity, with neither a principled notion of which items to forget nor any mechanism for distilling repeated episodes into reusable schemas. Systems like LangChain~\cite{langchain}, MemGPT~\cite{memgpt}, and Generative Agents~\cite{generative_agents} treat memory as a static or coarsely-paginated store, indiscriminately accumulating context until capacity limits are reached. As the store grows, retrieval precision deteriorates because the system lacks a mechanism to distinguish transient noise from persistent knowledge. This surfaces three concrete failure modes: (1) \textbf{memory bloat}, where token costs scale linearly with interaction length; (2) \textbf{semantic stagnation}, where specific episodes fail to evolve into generalised knowledge; and (3) \textbf{poor personalisation}, where a one-size-fits-all retention policy ignores individual user differences.

Recent work has addressed pieces of this problem in isolation but has not closed the loop. Generative Agents~\cite{generative_agents} introduce a reflection mechanism, but reflections are triggered by fixed rules and the per-tick re-processing cost is prohibitive online. Machine unlearning~\cite{sekhar2021unlearning} focuses on erasing data from model weights and is too mechanical to capture the salience-dependent, gradual forgetting seen in human memory. What is missing is a \emph{mechanism-level} framework that (i) forgets less aggressively when content is emotionally salient, (ii) consolidates repeated episodes into schemas without blocking user interaction, and (iii) adapts its own parameters to individual users.

To address this gap, we present \textbf{Mem0-Cognitive}, an \emph{emotion-weighted retention + asynchronous consolidation + adaptive reweighting} framework for LLM long-term memory. Rather than positioning our system as an instance of any particular cognitive-architectural theory, we describe it in mechanistic terms: three well-specified modules with code-level counterparts, each of which can be ablated independently. Specifically, the system (i) attaches an emotion intensity score $E_i \in [0, 1]$ to every stored memory $m_i$ and uses $E_i$ to stretch an Ebbinghaus-style retention curve (\S\ref{subsec:forgetting}); (ii) runs an offline consolidation cycle that clusters semantically-related memories and writes back consolidated records with a full audit log (\S\ref{subsec:consolidation}); and (iii) tunes per-user cognitive parameters through a top-$k$ reward-weighted-averaging heuristic (\S\ref{subsec:meta_learning}). The three modules are the only research claims this paper makes.

\textbf{Our Contributions} are:
\begin{enumerate}
    \item \textbf{A mechanism-level framework} for LLM memory management that combines emotion-weighted retention, asynchronous schema-style consolidation, and a top-$k$-weighted adaptive parameter tuner. Each module has a single public entry point in the accompanying code and can be enabled or disabled independently through a configuration flag.
    \item \textbf{An emotion-modulated retention law} (Eq.~\ref{eq:retention}) that couples an Ebbinghaus-style exponential decay with a real-time emotion-intensity term. The formula is monotonically non-decreasing in emotion at every fixed elapsed time and collapses to the classic Ebbinghaus curve when the emotion term is disabled; both properties are pinned by regression tests in the companion repository.
    \item \textbf{A planned benchmark}, \textbf{CognitiveBench}, whose seed-manifest schema and leak-check protocol are sketched in the appendix. The data generator itself and the ablation runner are explicitly marked as work-in-progress and are not shipped with this fork; a later version of this writeup would depend on implementing them.
    \item \textbf{An explicit alignment audit}. The companion repository is released with a documented mapping from every claim in this paper to the specific code path that implements it, including an explicit retraction of an earlier Gaussian-Process Bayesian-Optimisation claim (now described honestly as a top-$k$-weighted heuristic) and an explicit direction-correction of a retention-formula sign error present in an earlier version of the code.
\end{enumerate}

Consistent with this discipline, this version of the paper does \emph{not} report headline numeric results: the tables in \S\ref{sec:experiments} are intentionally empty until a generator and runner are implemented and emit stable outputs. We believe this \emph{artifact-first} stance is the honest position to take while the cognitive evaluation harness is still being built, and that explicit empty cells are a stronger signal of good-faith reproducibility than confidently-quoted numbers with unreachable scripts.
